\odiseachapter{Cuarta parte}{Calipso}

Definitivamente, los dioses son un problema.

Y lo son, por varias razones. La primera es que no paran de meter baza en los asuntos humanos, interviniendo, casi siempre para mal, tanto en la \emph{Ilíada} como en la \emph{Odisea}. El lector moderno no puede evitar sentirse frustrado por sus continuas injerencias, por la arbitrariedad de sus argumentos y, sobre todo, por su doble moral. Ya vimos, en el capítulo anterior, que toda la guerra de Troya podría haberse evitado si la diosa Discordia no hubiera sido tan artera, o las tres olímpicas principales se hubieran comportado como adultas, en lugar de mostrarse como unas malcriadas, acostumbradas a salirse con la suya.

Hay un ejemplo, en el libro IV de la \emph{Ilíada}, particularmente revelador. El canto abre con un concilio en el que los dioses debaten, mientras beben néctar, el destino de Troya. Zeus propone una salida razonable\footnote{La traducción de los versos que siguen es mía. Para no tener que compararme con Tobal, opto por los muy sufridos y castellanos endecasílabos.}:

\begin{verse}
¿O quizás reconciliar ambas partes?\\
Si la dulce paz nos conviene a todos,\\
salvarse puede la ciudad de Príamo,\\
y a Menelao le daríamos su esposa.
\end{verse}

Una solución la mar de razonable, que no complace ni a Hera ni a Atenea:

\begin{verse}
Sentadas una junto a otra, sus mentes,\\
fijas en arruinar a los troyanos.
\end{verse}

Hera, furiosa, arremete contra Zeus y le reprocha que después de haberse tomado ella tantos trabajos (para buscarles la desgracia a los de Ilión) quiera ahora el Crónida\footnote{Zeus es hijo de Cronos y, como buen hijo que es, le ha quitado el trono al padre y lo ha encerrado en lo más hondo del Tártaro.} resolverlo todo por las buenas. La salida de tono deja perplejo a Zeus, que le dice a su esposa (y hermana, lo del incesto tampoco se aplica a los olímpicos):

\begin{verse}
¡Oh, Hera, cuán extraña eres! ¿Qué\\
cuentas tienes con Príamo y sus hijos,\\
para siempre insistir en la ruina\\
de la ciudad de Troya de altos muros?

Solo te calmarías devorando\\
al buen Príamo y a todos sus vástagos;\\
devorarías a los troyanos todos,\\
para calmar tu implacable ira.

¡Pero sea! Por una menudencia\\
no permitamos brecha entre nosotros.\\
Atiende mis palabras sin embargo:\\
cuando sea yo quien destruir desee\\
una ciudad a la que tú favorezcas,\\
mi ira no intentes calmar por salvarles.
\end{verse}

Faltaría más que, por una ciudad más o menos, tuvieran que discutir los inmortales. La claridad con que Homero desvela su hipocresía, su falta de escrúpulos y su falta de compasión es apabullante.

Leer una historia en la que los dioses no hacen más que malmeter es frustrante, y escribirla todavía más. Por una parte, los olímpicos tienden a acaparar la atención más de la cuenta, desviándonos de las cuitas de los humanos, que son los que nos importan. Por otro lado, sus intervenciones destrozan cualquier argumento. El famoso \emph{deus ex machina}\footnote{La expresión significa literalmente «el dios desde la máquina»; en la práctica, el dios de la grúa. Cuando un dramaturgo se metía en un lío argumental del que no sabía salir, una tentación heredada de las épicas clásicas era resolverlo todo haciendo descender a un dios del cielo ---vía una grúa que sujetaba al actor de turno---.} aquí es literal. Dos guerreros, digamos, se enfrentan. El uno se ha preparado durante años para el combate. El otro es un golfo, pero sacrifica muchas reses a Zeus, o resulta ser uno de los cientos de hijos bastardos del Crónida. La historia pide a gritos que el guerrero que tanto tesón ha mostrado gane el combate, pero la mitad de las veces, según el humor del juntanubes, ocurre lo contrario. Los ejemplos no necesitan ser bélicos. Un hombre honesto puede verse robado por un tahúr que caiga bien a Atenea ---sin ir más lejos eso es lo que intenta Odiseo con los cícones--- o una bella joven se arriesga a ser mancillada por alguna deidad haciendo turismo sexual por la Tierra y, a poco que se descuide, represaliada por la esposa celosa de este. De ahí que muchas reescrituras modernas de los mitos griegos hayan optado por eliminar a los dioses o darles un papel marcadamente secundario.

En el caso del cine, el problema es todavía peor, porque cualquier puesta en escena de las deidades evoca en los espectadores las escenas de acción de mil películas de espada y brujería, o las versiones Marvel de esos mismos dioses.
Por ejemplo, el canto V de la \emph{Odisea} arranca con otro concilio de los dioses en el que Atenea vuelve a dar la lata ---lleva así desde el canto I--- con la suerte de Odiseo.

\begin{verse}
Ya de su lecho la Aurora, de junto a Titono adorable,\\
salía para traerles luz a sin muerte y mortales;\\
y se sentaban los dioses a junta, y Zeus tronante\\
en medio de ellos: poder es el suyo primero y más grande.\\
Y de Odiseo las cuitas Atena otra vez recordándoles\\
habla, pues daba en pensar que en poder de la ninfa él durase.
\end{verse}

Atenea quiere convencer a su padre de que ordene a la ninfa Calipso «soltar» a su protegido. Para situarnos: en los cuatro cantos (o libros) anteriores hemos seguido las andanzas de Telémaco, que pregunta primero a Néstor en Pilos y después a Menelao en Esparta por la suerte de su padre. Ulises todavía no ha aparecido en persona, aunque sabemos desde el arranque del poema que Calipso lo aloja como huésped un poco forzoso. Atenea dice así:

\begin{verse}
Zeus Padre y demás siemprevivas deidades benditas,\\
que en adelante no tenga un alma gentil, compasiva,\\
ya ningún rey, cetro en mano, ni piense en justicia divina;\\
no, que sea siempre intratable y actor de acciones impías,\\
y es que, ya veis, al divino Odiseo cualquiera ya olvida,\\
de esos que, siendo él señor, padre amable más bien parecía.
\end{verse}

¡La muy hipócrita! Apenas hace un ratito ---en el dilatado tiempo de los dioses--- no ha parado hasta arruinar Troya, importándole un bledo que Príamo fuera un rey justo o Héctor un héroe noble. De hecho, como Nolan muestra en la escena del caballo, los troyanos rinden culto a Atenea, que no solo tiene un templo en el alcázar, sino que ostenta el título de \gr{ἐρυσίπτολις} (erysíptolis), «protectora de la ciudad». ¡Valiente protectora! No ceja hasta que arde Troya y los fieles que le han ofrecido todo tipo de sacrificios son pasados por el cuchillo. Pero cuando se trata de Ulises, no tiene problema en invocar la justicia divina que ella se salta a la torera. Su discurso sigue así:


\begin{verse}
Llevando recios dolores, yaciendo ahora está en una isla,\\
en la mansión de la ninfa Calipso, que allá lo esclaviza\\
bien por la fuerza; y su tierra llegar a alcanzar no podría,\\
pues que ni naves remeras él tiene ni feligresía\\
que lo llevara hasta allí por las anchas espaldas marinas.\\
Y ahora a su hijo querido por darle muerte suspiran\\
de vuelta a casa; marchó, por si allí de su padre algo oía,\\
a la santísima Pilos y a Lacedemonia divina.
\end{verse}

%¿Es cierto que Calipso lo esclaviza, «bien por la fuerza»? Ya vimos en el capítulo anterior que Helena era una narradora de la que no podíamos fiarnos y es bien posible que a Atenea le pase lo mismo.

Zeus empieza por contestarle que si tanto aprecia a Odiseo, ponga manos a la obra para ayudarlo.

\begin{verse}
Y las palabras de vuelta que Zeus juntanubes le daba:\\
«¡Hija, qué mala palabra saltó de tus dientes la tapia!\\
¿No te has resuelto tú misma contigo sola esa traza\\
de que de aquellos se vengue Odiseo al volver a su casa?\\
Guía a Telémaco tú, que es tarea que puedes muy sabia,\\
para que, salvo y sin mal, alcance su tierra y su patria;\\
y pretendientes y nave, que como llegaron se vayan».
\end{verse}

Pero también decide tomar cartas en el asunto:

\begin{verse}
Decía; y ahora es a Hermes, su hijo, al que mira y ordena:\\
«Hermes, pues que eres tú siempre el que lleva recados en estas,\\
nuestro infalible querer ve a decirlo a esa ninfa beltrenza:\\
que es de Odiseo biensufridor ya el viaje de vuelta\\
sin que de dioses escolta ni de morideros le venga;\\
que en balsa milbientrabada, pasando él solo miserias,\\
en veinte días podría llegar a la fértil Esqueria,\\
tierra que lo es de feacios, que de los dioses van cerca\\
y honra de dios le darían aquellos en sus entretelas,\\
y en un navío a su patria lo habrán de mandar y a su tierra,\\
luego de darle vestidos y bronce y oro a lleneza,\\
mucho, que nunca Odiseo de Troya tanto trajera\\
si ileso hubiera venido, escogida su parte de presa.\\
Ver a los suyos, llegar a su casa techalta y su tierra\\
cosa es que ya bien dispuesta la tiene la providencia».
\end{verse}

Relee con calma, biensabida lectora, todoatento lector, los versos y disfruta con «beltrenza», «biensufridor», «milbientrabada» y «techalta», pero repara también en que Zeus tiene claro todo el argumento, dejándole muy poco juego al azar. Calipso debe soltar a Ulises, que en veinte días ---ni más ni menos--- va a llegar a la isla de los feacios donde será colmado de regalos y devuelto a su casa. Es decir, cuando se pone, Zeus sabe ---o decide--- exactamente qué pasa. De ahí que los diez años que Ulises lleva dando tumbos se nos antojen innecesarios. A estas alturas ya le ha pasado de todo. Se ha encontrado con Polifemo y los lestrigones, ha escuchado a las sirenas, ha pasado entre Escila y Caribdis, se ha parado un año en casa de Circe ---de la que hablaremos pronto---, ha bajado a los infiernos y regresado, ha naufragado, perdiendo a todos sus compañeros. Pero nada de eso hacía falta. Si a Zeus le hubiera apetecido, habría llegado a casa en quince días. Hay un elemento de arbitrariedad insoportable en estos dioses, que lo mismo te dan la mano que te tiran un rayo, que lo mismo se te aparecen a darte conversación, que se olvidan de que existes durante siete años. No es de extrañar que Alessandro Baricco, por ejemplo, los quite del medio en su solvente reescritura de la \emph{Ilíada} o que Nolan intente, desesperadamente, que no le den mucho la lata. Imagina, cinéfila lectora, fílmico lector, que hubiera tenido que rodar la siguiente escena:

\begin{verse}
Tal dijo, y no desoyó el correo Argifontes la manda.\\
Conque enseguida a sus pies se amarró las preciosas sandalias\\
de oro, no humanas, que a él se lo llevan por sobre las aguas\\
y por la tierra sin linde, igual que de un viento una ráfaga.\\
Tomó en su mano la vara con que a hombres ensueña y encanta,\\
siempre que quiere, los ojos, y a los que se aduermen espanta.\\
Ella en la mano, el pujante Argifontes de vuelo marchaba.\\
Atravesando Pieria, del cielo se echó a la maralta;\\
y se lanzó a flor de ola como una gaviota que, caza\\
dándole al pez por las ollas tremendas del mar sin segada,\\
moja sus alas espesas por las milespumas saladas:\\
a ella semejo, las olas a Hermes llevábanlo, tantas.
\end{verse}

¡Hermes, con cetro y sandalias aladas volando sobre las aguas! ¡A ver quién se atreve en una película que quiere ser seria sin que el público vea al Flash de DC, con uniforme y todo!

\begin{verse}
Cuando por fin arribó a aquella isla remota,\\
se levantó de las aguas violeta y, hollando la costa,\\
fue caminando hasta dar en la grande caverna en que mora\\
la ninfa: dentro ella estaba y allí la encontró, rubihermosa.\\
En el fogón crepitaba un gran fuego, y se olía el aroma\\
del cedro bienastillado y el limonero por toda\\
la isla en su arder; ella dentro, cantando con voz hechizosa,\\
puesta al telar, lanzadera de oro, tejía allí sola.
\end{verse}

No es una gruta cualquiera, la de Calipso rubihermosa. Más bien un hotel de cinco estrellas:

\begin{verse}
Exuberante, a redor de la gruta verdeaba una fosca\\
de aliso y chopo y ciprés bienoliente, bajo cuyas hojas\\
iban a hacerse sus nidos las aves alidonosas,\\
los búhos y los falcones y las lengüilárguidas chovas\\
marinas, que ellas su jera la tienen entre las olas.\\
Y allí también prosperaba, de la honda cueva a redonda,\\
viña labrada en sazón cargada de pámpanos toda.\\
Cuatro fontanas hacían con su agua clara una horca\\
corriendo a parte y a parte, al pie cada cual de la otra.\\
De un lado y otro mullidos de perejil y de viola\\
se florecían los prados. Llegar hasta allí y ver es cosa\\
que en sus adentros hasta a uno sin muerte lo atrapa y emboba.
\end{verse}

Si incluso un inmortal se queda embobado con las delicias del lugar, ¿debemos creernos que Ulises lo está pasando fatal?

Pues el caso es que así es.

\begin{verse}
Pero a Odiseo grancorazón allí dentro no encuentra,\\
no, que lloraba en el alto cantil ---¡cuántas veces lo hiciera!---\\
dilacerándose el alma en lágrimas, plantos y penas.\\
Ve el mar, el mar malagosto, y correr sus lágrimas deja.
\end{verse}

Calipso, por su parte, se huele que la visita no es de simple cortesía:

\begin{verse}
Y preguntando está a Hermes Calipso, la dea de deas,\\
luego de hacerlo sentar en asiento esmaltado de estrellas:

«¿A qué te me eres venido, Hermes varitadeoro\\
tú, tan honrado y querido? Hasta hoy te me dabas bien poco.\\
Di lo que quiera que pienses, que está, sí, en cumplirlo mi antojo,\\
si es que lo puedo cumplir y si es cumplidero el propósito.\\
Ah, pero sígueme adentro, que honores de huésped dispongo».
\end{verse}

Hermes acepta la hospitalidad y, después de comer y beber a gusto, entra en materia. Órdenes de Zeus, le comunica:

\begin{verse}
Dice que un hombre hay contigo, el más infeliz de la copia\\
de hombres que en torno a los muros, por nueve años, de Troya\\
lucharon, y que, asolando al décimo aquella almodóvar,\\
se regresaron; mas, ¡ay!, que ofendieron a Atena en su torna,\\
y levantó para ellos mal viento y altas las olas.\\
Todos los otros murieron entonces ---su válida tropa---,\\
pero a él el viento y las olas lo echaron aquí. A ti ahora\\
te manda que lo despidas, y sin ninguna demora.\\
No es acabar aquí lejos, sin deudos, lo que le toca;\\
no, que pisar tierra patria y la casa techalta ya es cosa\\
que, a más de ver a los suyos, dejó la parca bien pronta.
\end{verse}

Las órdenes, a la ninfa, le sientan fatal:

\begin{verse}
Tristes sois, dioses, que andáis ---y a lo más--- envidiando a los otros.\\
Y os enrabiáis con las diosas que a hombres, y sin rebozo,\\
llevan al lecho, si una al que quiere lo tiene de esposo.\\
Cuando raptó a Orión la Aurora dedos de rosa a vosotros,\\
dioses de fácil vivir, os vino a entrar tal enojo,\\
que allá en Ortigia la muerte casta Ártemis tronodeoro\\
le hubo de dar con sus flechas de dulce pasión. De igual modo\\
cuando a Jasión la belrubia Deméter, rendida a su antojo,\\
se unió en amor y delicias sobre un barbecho esponjoso\\
requetearado, no mucho tardó en saber eso todo\\
Zeus, y blanca centella soltó, y lo mató en ese pronto.
\end{verse}

La rabieta de Calipso tiene una lectura interesante. Según ella, no se trata de que los olímpicos se preocupen por el destino de un hombre como Ulises, sino que les da rabia que las diosas se encamen con simples mortales. Y sigue:

\begin{verse}
Así conmigo ahora, dioses, os enrabiáis porque gozo\\
de hombre mortal que salvé, montado aquel sobre un tronco,\\
solo, cuando con centella de plata su barco afanoso\\
milastilló Zeus en medio del vinoscuro marhondo.
\end{verse}


Como ya sabíamos. El mismo dios que pretende ahora devolverlo a casa le hizo naufragar por un quítame allá estas vacas.

\begin{verse}
Todos los otros murieron entonces, sus buenos devotos,\\
pero a él el viento y las olas aquí lo acercaron. Yo solo\\
era en amarlo, en cuidarlo, en dar en decir cada poco\\
que yo lo haría inmortal, sin vejez para ya el tiempo todo.
\end{verse}

Y aquí, en estas precisas líneas, Calipso se gana a cualquiera. Lo salva, lo cuida, lo mima, le ofrece la inmortalidad. No cabe duda de que el héroe le gusta, pero hay algo más: Calipso está sola. La película de Nolan lo captura implícitamente, entre otras cosas gracias a la acertada interpretación de la divinal Charlize Theron. Ninfa y todo, no tiene quien la quiera, vive abandonada en su preciosa isla, en su cómoda gruta, sin que nadie se preocupe por ella. ¿Cómo no comprenderla?

\begin{verse}
Mas si el querer de Zeus cabrillante no hay forma ni modo\\
de que otro dios lo rebase ni de que lo deje sin logro,\\
váyase, si es que de aquel se me vienen mandado y exhorto,\\
por ese mar malagosto: yo hacerlo partir no sé cómo,\\
pues ni con naves remeras yo cuento ni con los devotos\\
que por el ancho lomo del mar con él fueran de acólitos.\\
Aunque consejos no me ahorraré ---corazón aquí pongo---\\
para que, salvo y sin mal, se le dé a esa su tierra el retorno.
\end{verse}

La ninfa, qué remedio, acata órdenes y a continuación dice que no quiere saber nada de la operación rescate. Hermes, discretamente, la amenaza:

\begin{verse}
Y le decía aquí a ella a su vez el correo Argifontes:\\
«Hazlo partir así ya, que la ira de Zeus no te toque,\\
no vaya a ser que te coja rencor y se ensañe a mayores».
\end{verse}

Hermes se va y la siguiente escena revela otro dato interesante:

\begin{verse}
A la que dijo él así, partió el Argifontes pujante.\\
La ninfa amante a buscar a Odiseo corazongrande\\
fue, pues que bien escuchó y dio oído de Zeus al mensaje.\\
Y lo encontró sobre el alto cantil: de sus ojos secarse\\
nunca su llanto dejaba; rodaba su edad dulceamable en\\
lágrimas por su regreso; la ninfa dejó de gustarle.
\end{verse}

«Dejó de gustarle», por lo que podemos suponer que alguna vez sí le gustó, pero el caso es que ahora está harto:

\begin{verse}
Que es que ---¿qué cosa iba a hacer?--- lleva noche tras noche acostándose\\
él sin deseo en las cuéncavas grutas con ella deseándole.\\
Y las mañanas, sentado entre peñas y por arenales\\
y desgarrándose el alma en sollozos, quejumbres y males,\\
ve solo el mar, ese mar malamiés que las lágrimas trae.
\end{verse}

Se acuesta sin deseo (ella sí lo tiene) y se muere de nostalgia. ¿Desde hace cuánto tiempo? Uno de los «misterios» de la \emph{Odisea} es por qué Ulises se demora tanto donde Calipso. La explicación más natural es que la bellitrenza no está dispuesta a dejarlo partir hasta que llegan órdenes de la superioridad. En su versión, Nolan opta por otra explicación. Ulises, atormentado por el sentimiento de culpa, se convierte en un yonqui y la ninfa no tiene problema en proporcionarle su droga cada noche. La «explicación» encaja con el argumento general del director, que hace girar la personalidad del héroe en torno al arrepentimiento por sus acciones en Troya ---convirtiéndolo así en un personaje bastante plano---. El recurso no está mal pensado ni queda tan lejos de la lectura literal, en la que la ninfa lo retiene a la fuerza. En la película también lo hace, pero a base de drogas. Por otra parte, cuando ya no le queda otra, Calipso muestra una ternura que no es corriente entre los inmortales ---y que Charlize Theron escenifica muy bien---:

\begin{verse}
«Ya, tristestrella, no más ya me llores aquí, que tu tiempo\\
no se te vaya; que ya te despido, que ya estoy en ello.\\
Arma con bronce una balsa cortando unos largos maderos,\\
de buena anchura, y clava un castillo bien alto allí luego\\
para que pueda llevarte por el hondomar neblinegro.\\
Yo pondré el pan, pondré el agua, el rojo vino a contento\\
yo lo pondré, porque el hambre de ti se tenga bien lejos;\\
y te daré unos vestidos, y te mandaré buen oreo,\\
para que, salvo y sin mal, se te dé a esa tu patria el regreso,\\
si así lo quieren los dioses que tienen el ancho cielo,\\
ellos, más altos que yo en saberes y cumplimientos».
\end{verse}

Resignada, pero generosa. Bueno, no del todo resignada. Un poco más adelante vuelve a jugar su mejor carta:

\begin{verse}
«Sangre Laercia, Odiseo milmañas, hijo del cielo,\\
¿entonces ahora a tu casa, a tu patria querida, a tu pueblo\\
quieres ya mismo volver? Buena suerte y adiós, yo obedezco.\\
Mas si supiera tu corazón la ración de los duelos\\
que hadada está para ti, hasta dar en el suelo paterno,\\
aquí conmigo quedaras a ser guardador de este techo,\\
y un inmortal, por muy mucho que ahora te tenga ese anhelo:\\
ver a tu esposa, que así es cada día y en todo momento.\\
Y no, no soy inferior yo a ella ---gloriarme bien puedo---\\
en porte ni en condición, y mujer moridera no es bueno\\
que entre a porfía con diosas por cuál estatura y qué cuerpo».
\end{verse}

El soborno no es pequeño: la vida eterna, nada menos. A la ninfa, como es natural, no le acaba de entrar en la cabeza que Odiseo prefiera irse a buscar a una mujer mortal que, con certeza, no puede compararse con su divina presencia. Y aquí es donde vuelve a verse el miedo de Hollywood a Homero. Una cosa es el estrés postraumático, más fácil de digerir que un yogur. Y otra es darle a Ulises todas sus dimensiones, que van desde la crueldad más extrema ---no hay pista alguna en la \emph{Odisea} de que se arrepintiera de sus andanzas en Troya, y la masacre de los pretendientes y las criadas deja claro que no anda muy sobrado de compasión--- hasta la abnegación más absoluta ---renunciar a la vida eterna a cambio de volver a casa---. Ciertamente no es fácil y, de hecho, la \emph{Odisea} no se mete en camisa de once varas: se limita a mostrarnos todas las facetas del héroe sin pretender reducirlo a un monotema, que es, en mi opinión, el fallo fatal de la película de Nolan.

Y eso, a pesar de que el tema del arrepentimiento es muy fértil. En mi \emph{Abandonando Ítaca}\footnote{\url{https://eolasediciones.es/libro/abandonando-itaca-leaving-ithaca/}} juega un papel central. Pero el truco que me permite abordarlo sin meterme en líos es que \emph{Abandonando Ítaca} arranca veinte años después de que Ulises llegara a la isla, cuando él es ya un anciano y casi todo el mundo que conoce ---incluyendo Penélope--- ha muerto.


\begin{verse}
I

Inesperadamente, has llegado a esa edad en la que los fantasmas\\
superan en número a los vivos. Están por todas partes,\\
acechando entre sombras, todos esos amados\\
y muertos rostros, disolviéndose en la niebla del tiempo.

II

Rostros muertos. Anticlea, tu madre, murió en Ítaca,\\
mientras tú estabas ocupado, masacrando a los hijos\\
de las mujeres de Troya. Y, sin embargo,\\
pareció tan injusto volver tan tarde para algo que no fuese\\
quemar sobre su tumba incienso, manchada tu cabeza con cenizas.

III

Y Laertes más tarde. Te esperó,\\
y junto a ti celebró la carnicería de los jóvenes,\\
pero murió poco después. Ahora acude a tus sueños\\
a menudo, como si quisiera decirte algo esencial,\\
sobre la vida y la muerte. Algo que siempre,\\
simplemente escapa de tu comprensión.\\
Tú, que has matado a tantos hombres,\\
sin preguntar nunca por qué.

Sin embargo, tu padre era un hombre discreto,\\
y su fantasma no olvida sus modales.\\
No se demora demasiado, y su presencia no es fatigosa nunca.

IV

Tampoco la muerte ha cambiado a Penélope;\\
su fantasma es tan obstinado como la mujer\\
que te estuvo esperando veinte años.\\
Cada día tejía su mortaja, solo para\\
destejerla a su vez a la llegada de la noche.\\
Semanas que se convertían en meses y meses que se convertían en años,\\
y así a lo largo de dos décadas. ¿Entonces,\\
te sorprende que siga viniendo, antes del alba,\\
y se siente a tu lado para usar sus agujas?

No te importaría tanto si de vez en vez,\\
ella dijera algo.\\
Quedó tanto sin decir entre vosotros\\
que tú nunca te atreviste a contar,\\
y que ella nunca preguntó.\\
De todos modos, ¿de qué habría servido?\\
Nadie intenta explicar las andanzas de un héroe,\\
y ninguna esposa razonable\\
esperaría la castidad que, en cambio, se daba por supuesto\\
que ella debía observar.

Ahora se ha ido. Murió tan en silencio\\
como vivió. Mas su fantasma no se irá,\\
has intentado explicar tus razones y justificarte,\\
todo en vano. Ella, paciente y muda,\\
volverá cada noche, simplemente\\
para tejer y destejer.
\end{verse}


Es una entrada bastante distinta de la de Nolan, sobre todo en lo que se refiere a Penélope, un personaje bastante más ambiguo de lo que el director nos muestra en la película, donde solo la pasión de Anne Hathaway consigue salvarla ---por los pelos--- del cliché. En mi versión, Penélope cumple, obstinadamente, como en la \emph{Odisea}, pero sus razones se le escapan al héroe y al lector.

No es de extrañar que, anciano, viudo, solitario, a las puertas de la muerte, Ulises reflexione sobre el pasado:

\begin{verse}
V

Rostros muertos. Cuentan incluso los de tus enemigos,\\
te aparecen en sueños, no enfadados, tristes,\\
como si preguntasen. ¿Crees que fue necesario?\\
¿Era realmente imperioso lanzar desde los muros\\
de la ciudad en ruinas a ese pobre bebé? Pero en aquel momento,\\
tu corazón estaba tan lleno de rabia y tú tan orgulloso\\
de tu propia crueldad. Ahora el rostro\\
del niño muerto vive en tus pesadillas. Está ahí, y junto a él,\\
los soldados troyanos, sacrificados como reses,\\
para aplacar la ira de los aqueos\\
y entretener el capricho de los dioses.

VI

Pero lo más vívido de todo es ver los rostros\\
de esos locos muchachos, esa alegre pandilla de borrachos,\\
que tan solo querían divertirse a tu costa.\\
Es cierto, se comieron tus corderos,\\
se bebieron tu vino, pretendiendo cortejar a tu mujer,\\
aunque en realidad no le tocaron ni un pelo.\\
¿Por qué tanta violencia entonces, por qué la carnicería?\\
¿En el nombre de qué? ¿Del honor? ¿Por venganza?\\
Dime: ¿realmente lo pasaste tan mal en casa de Calipso?\\
¿Negarás que el amor de Circe era tan ardiente\\
como el mismísimo infierno que solo por ella recorriste?
\end{verse}

¿Por qué la carnicería? No hay versión de la \emph{Odisea} que no se tropiece con ella y todas intentan justificarla como necesaria, cuando no lo era en absoluto ---y mucho menos la ejecución posterior de las esclavas---. Pero es que Hollywood teme la crueldad de Odiseo y sus tremendos defectos, que lo hacen muy poco edificante como modelo a seguir. Y, sin embargo, Homero los muestra sin recato alguno.

Ocurre, además, que no siempre es cruel ni vengativo. Con Calipso se porta como un caballero. Al discurso de la bellitrenza, responde:


\begin{verse}
Y el milardido Odiseo así respondiéndole dijo:\\
«No te me aíres por eso, diosa señora, yo mismo\\
sé que Penélope milpensamientos de ti está lejísimos\\
en estatura y belleza, no hay más que mirar lo que digo.\\
Y ella es mortal, y la muerte y vejez nada tienen contigo.\\
Pero, con todo, yo quiero, y todos los días lo pido,\\
ver día en que uno regresa, y a casa llegar. Y si vivo\\
nuevo revés por el mar vinoscuro de algún luzdivino,\\
aguantará el corazón en mi pecho, en dolores curtido;\\
sí, que mucho es lo sufrido y es mucho lo que he resistido\\
por mar y en guerra: que en pos de esos tantos sea el nuevo venido».
\end{verse}

Odiseo milmañas no es siempre un personaje simpático, ya lo hemos dicho. Pero aquí demuestra a la vez ternura hacia la desilusionada ninfa, diplomacia y valor. Ni quiere ser inmortal ni le importa seguir jugándosela con tal de llegar a casa. Pero lo cortés no quita lo valiente (o más bien lo contrario en este caso). Una vez resuelta la situación, se impone una despedida «como dios manda» (esa era una frase favorita de mi padre) o, más bien, «como la diosa quiere».

\begin{verse}
Dijo. Y el sol se ponía, y venía cayendo la sombra.\\
Y entrando al seno profundo los dos de la gruta milhonda\\
gozan amor y delicias el uno enredado en la otra.
\end{verse}
